\section{结论}\label{sec:conclusion}

本文研究闭卷条件下如何利用 LLM 的分析来检测虚假新闻。我们提出 \method{}：LLM 把新闻分解为声明并为每条声明写出正反两方论证，可训练的裁决器以声明为实例做多示例学习，用相关性门控剔除不属于该新闻的分析单元，并在所有单元都不可用时退回到只依赖文本的判断。我们在统一了 LLM、超参数搜索与随机种子的条件下，把它与 ARG、FactGuard、直接微调的 LLM 以及一个简单的评分融合基线做了比较，并设计了四项受控测试。

结论有三层。第一，就方法而言，\method{} 在分布内与已有的协同方法持平；它的相关性门控是有效的，能可靠地识别与新闻对不上的分析单元，但这回答的只是“分析是否属于这条新闻”，而不是“分析是否正确”。第二，就这一类方法而言，我们没有找到能在闭卷条件下估计“LLM 的判断是否可靠”的信号：\method{} 的判定门控与 ARG 的可用性机制实际学到的都是对新闻标签的预测。第三，就任务而言，风格改写与事实篡改两项测试表明，包括 LLM 自身判断在内的所有闭卷检测器主要依赖文风，对内容层面的事实改动几乎没有反应；把 LLM 的分析做得更细、更结构化，并不能改变它所依据的信息。

这些结果指向两个后续方向。其一，要让检测器真正对“内容是否属实”作出反应，需要引入模型参数之外的信息，例如检索到的证据；本文的声明分解与相关性门控可以直接用于判断检索到的证据是否与某条声明相关。其二，评测上，闭卷检测方法应当在风格受控与事实受控的测试集上报告结果，并与“直接微调同一个 LLM”以及“融合一个 LLM 评分”这两个基线比较；在多随机种子的设置下，1 个点以内的差距不足以支撑结论。
